\documentclass[12pt, a4paper]{ctexart}
\usepackage[margin=2cm]{geometry}
\usepackage{libertine}

\usepackage{titlesec}
\usepackage{zhnumber}
\titleformat*{\section}{\Large\bfseries\raggedright}
\renewcommand{\thesection}{\zhnum{section}}
\renewcommand{\thesubsection}{\arabic{subsection}}

\setcounter{secnumdepth}{4}
\renewcommand{\theparagraph}{\thesubsubsection.\arabic{paragraph}}
\renewcommand{\paragraph}[1]{
    \refstepcounter{paragraph}
    {\bfseries\theparagraph\quad#1}\par
    \vspace{2pt}
    \noindent
}

\usepackage{enumitem}
\setlist[enumerate]{itemsep=2pt, parsep=0pt, partopsep=0pt, topsep=2pt}
\setlist[itemize]{itemsep=2pt, parsep=0pt, partopsep=0pt, topsep=2pt}
\linespread{1.2}

\usepackage{minted}
\setminted{
    breaklines=true, escapeinside=||, fontsize=\small, frame=lines,
    mathescape=$$, numbers=none, style=bw, tabsize=4
}

\usepackage{graphicx}
\usepackage{siunitx}

\begin{document}

\pagestyle{plain}
\thispagestyle{empty}

\noindent
\begin{tabular*}{\textwidth}{l @{\extracolsep{\fill}} r @{\extracolsep{6pt}} l}
    \LARGE{\textbf{实验报告}} & 数据结构 & \textit{Data Structures} \\
\end{tabular*}\\\\
\begin{tabular*}{\textwidth}{l l}
    \textbf{报告标题: } & \textbf{二叉排序树} \\
    \textbf{学号: } & 19240212 \\
    \textbf{姓名: } & 华博文 \\
    \textbf{日期: } & 2025 年 12 月 16 日 \\
\end{tabular*}\\
\rule[2ex]{\textwidth}{2pt}

\section{实验环境}

\subsection{操作系统}

Ubuntu 24.04.3 LTS x86\_64，Linux 6.14.0-36-generic，AMD Ryzen 7 7700 (16) @ 5.3GHz，DDR5。

\subsection{编程工具}

NeoVim v0.12.0-dev，gcc 13.3.0。

\subsection{其他工具}

\LaTeX{} (\TeX\textit{studio})。

\section{实验目的}

\begin{enumerate}
    \item 理解二叉排序树（BST）的定义、性质，与线性表、一般二叉树的区别；
    \item 掌握 BST 的插入、查找（递归 / 非递归）与中序遍历的实现，理解中序与有序性的关系；
    \item 学习平均查找长度（ASL）的计算方法，分析树形态对查找性能的影响，比较退化链表与平衡情形的复杂度差异。
\end{enumerate}

\section{实验内容及其完成情况}

\begin{enumerate}
    \item 从无序整数序列构建 BST：设计节点结构，按输入顺序插入，验证中序遍历为升序；
    \item 查找操作：实现迭代查找并返回路径，同时暴露指向存储元素的指针；
    \item 性能指标：通过一次遍历统计各节点深度，计算 ASL；
    \item 输入 / 输出：支持从文件或标准输入读取，交互式查询任意键值，异常输入报错。
\end{enumerate}

\subsection{数据结构}

\begin{itemize}
    \item 模板 BST：\mintinline{c++}`template <typename T> class bst`，节点含键值与左右子指针；
    \item 查找结果：\mintinline{c++}`find_result`，包含 \mintinline{c++}`ptr`（指向存储元素，未命中为空）与比较路径 \mintinline{c++}`path`；
    \item 遍历辅助：\mintinline{c++}`traverse(Fn &&)`，向回调传递键与深度，用于一次遍历计算 ASL。
\end{itemize}

\subsection{程序实现}

首行整数 $n$，其后 $n$ 个整数（用空格/换行分隔）。传入 \mintinline{text}`-` 时从标准输入读取。

\begin{minted}{c++}
template <typename T>
class bst {
public:
    struct find_result { const T *ptr{nullptr}; std::vector<T> path; };
    void insert(const T &key) { insert(root_, key); }
    void build(const std::vector<T> &keys) { for (const T &v : keys) insert(v); }
    std::vector<T> inorder() const { std::vector<T> out; inorder(root_, out); return out; }
    find_result find(const T &key) const {
        std::vector<T> path; const node *cur = root_.get();
        while (cur) {
            path.push_back(cur->key);
            if (key == cur->key) return {&cur->key, std::move(path)};
            cur = (key < cur->key) ? cur->left.get() : cur->right.get();
        }
        return {nullptr, std::move(path)};
    }
    bool empty() const { return !root_; }
    template <typename Fn>
    void traverse(Fn &&fn) const { traverse(root_.get(), 1, std::forward<Fn>(fn)); }
};

template <typename T>
double compute_avg_path(const bst<T> &tree) {
    if (tree.empty()) return 0.0;
    long long sum = 0, cnt = 0;
    tree.traverse([&](const T &, int depth) { sum += depth; ++cnt; });
    return cnt ? static_cast<double>(sum) / cnt : 0.0;
}

std::vector<int> load_data(const std::string& path);

int main(int argc, char **argv) {
    std::vector<int> data = load_data(argv[1]);
    bst<int> tree; tree.build(data);
}
\end{minted}

\subsection{编译与运行示例}

在 \mintinline{text}`src` 目录下编译：

\begin{minted}{bash}
g++ -std=c++17 -O2 main.cpp -o bst
\end{minted}

使用示例数据 \mintinline{text}`res/sample.txt`，从文件读取后可继续交互查询：

\begin{minted}{bash}
./bst ../res/sample.txt
\end{minted}

或使用标准输入：

\begin{minted}{bash}
cat ../res/sample.txt | ./bst -
\end{minted}

\subsection{样例输入输出}

输入文件 \mintinline{text}`res/sample.txt`：

\inputminted{text}{../res/sample.txt}

运行示例（节选）：

\begin{minted}{text}
In-order traversal (ascending): 10 12 23 34 45 50 56 67 78 89
Average successful search length (ASL): 2.9
Search key (q to quit): 50
Path: 45 -> 78 -> 56 -> 50
Found key 50
Search key (q to quit): 13
Path: 45 -> 12 -> 34 -> 23
Key 13 not found
Search key (q to quit): q
\end{minted}

\subsection{复杂度分析}

\begin{itemize}
    \item 构建：插入 $n$ 个元素，时间 $\mathcal{O}\left(n\cdot h\right)$，最坏退化 $\mathcal{O}\left(n^2\right)$，平均 / 平衡约 $\mathcal{O}\left(n\log n\right)$；空间 $\mathcal{O}\left(n\right)$；
    \item 查找：时间 $\mathcal{O}\left(h\right)$，最坏 $\mathcal{O}\left(n\right)$，平衡 $\mathcal{O}\left(\log n\right)$；
    \item ASL 统计：一次遍历 $\mathcal{O}\left(n\right)$ 时间，$\mathcal{O}\left(1\right)$ 额外空间（递归深度按树高计）。
\end{itemize}

\section{结论}

完成了 BST 的构建、迭代查找与路径输出，基于一次遍历计算 ASL，支持文件/标准输入与交互查询。实验验证中序遍历有序性，演示了查找路径与命中情况，并分析了平衡与退化形态下的复杂度差异。

\end{document}
